\chapter{Maximum Flow and Matching}
\chapterepigraph{Maximum flow is the master problem of this chapter. Once
you can push the most units from a source to a sink through capacitated
edges, you get minimum cuts, bipartite matchings, and edge-disjoint paths
almost for free -- each is a maximum flow with the graph set up just so.}

\section{Everything Is a Flow}
\label{sec:cses-graph-flows-intro}

A flow network has capacities on edges; the maximum flow from source $s$
to sink $t$ is the greatest total that can be pushed while respecting
capacities and conservation. Augmenting-path algorithms (BFS-based
Edmonds--Karp, or Dinic with level graphs) compute it.

\begin{patternbox}[Flow Reductions in This Chapter]
\begin{itemize}
  \item \textbf{Maximum throughput / edge-disjoint paths count} --
        \textbf{max flow} with unit or given capacities: Download Speed,
        Distinct Routes (unit capacities; then decompose the flow into
        paths).
  \item \textbf{Fewest edges to disconnect $s$ from $t$} -- \textbf{min
        cut} $=$ max flow (max-flow--min-cut theorem), and the cut edges
        are those from reachable to unreachable after saturation: Police
        Chase.
  \item \textbf{Maximum bipartite matching} -- \textbf{max flow} with a
        super-source into one side and a super-sink out of the other, all
        unit capacities: School Dance.
  \end{itemize}
\end{patternbox}

\begin{ideabox}[Residual Graphs and Augmenting Paths]
Flow algorithms repeatedly find a path from $s$ to $t$ with spare capacity
and push along it, adding \emph{reverse} edges so earlier decisions can be
undone. When no augmenting path remains, the flow is maximum, and the
nodes still reachable from $s$ define the minimum cut.
\end{ideabox}

\newpage

\section{Download Speed}
\label{sec:cses-download-speed}

\problemstatement{Given a directed graph of $n$ nodes and $m$ edges, each
with a capacity (bandwidth), find the maximum data rate from node $1$
(source) to node $n$ (sink).}

\begin{patternbox}
Maximum throughput under capacities is the \textbf{maximum flow} problem.
\textbf{Dinic's algorithm} computes it efficiently: repeatedly build a
level graph with BFS, then push blocking flow along shortest augmenting
paths with DFS.
\end{patternbox}

\subsection*{Level graphs and blocking flow}

A flow assigns each edge a value up to its capacity, with flow conserved at
every node except source and sink; the max flow is the greatest total
leaving the source. Dinic alternates two phases: a BFS labels each node with
its distance (level) from the source in the residual graph; a DFS then sends
as much flow as possible along edges that go strictly one level deeper (the
blocking flow). Reverse edges are maintained so earlier flow can be
rerouted. When BFS can no longer reach the sink, the flow is maximum.

\begin{ideabox}[Augment Along Shortest Paths, in Bulk]
Restricting augmenting paths to the level graph (only forward-level edges)
and saturating them all before rebuilding gives Dinic its speed: the number
of BFS phases is $O(V)$, far fewer than pushing one path at a time.
\end{ideabox}

\subsection*{Algorithm}

\begin{enumerate}
  \item Build the flow network with forward capacities and zero-capacity
        reverse edges.
  \item While BFS reaches the sink: DFS blocking flow along level edges,
        accumulating flow.
  \item Output the total flow.
\end{enumerate}

\subsection*{Dry run}

\dryrun{$1\!\to\!2(3)$, $1\!\to\!3(2)$, $2\!\to\!4(2)$, $3\!\to\!4(3)$;
sink $4$.}

\begin{itemize}
  \item Push $2$ along $1\!\to\!2\!\to\!4$ and $2$ along
        $1\!\to\!3\!\to\!4$.
  \item Max flow $=4$ (limited by outgoing/incoming capacities).
\end{itemize}

\subsection*{C++ implementation}

\begin{lstlisting}
#include <bits/stdc++.h>
using namespace std;

struct Dinic {
    struct Edge { int to; long long cap; int rev; };
    vector<vector<Edge>> g;
    vector<int> level, it;
    Dinic(int n) : g(n), level(n), it(n) {}
    void addEdge(int u, int v, long long c) {
        g[u].push_back({v, c, (int)g[v].size()});
        g[v].push_back({u, 0, (int)g[u].size() - 1});   // reverse edge
    }
    bool bfs(int s, int t) {
        fill(level.begin(), level.end(), -1);
        queue<int> q; level[s] = 0; q.push(s);
        while (!q.empty()) {
            int u = q.front(); q.pop();
            for (auto& e : g[u])
                if (e.cap > 0 && level[e.to] < 0) { level[e.to] = level[u] + 1; q.push(e.to); }
        }
        return level[t] >= 0;
    }
    long long dfs(int u, int t, long long f) {
        if (u == t) return f;
        for (int& i = it[u]; i < (int)g[u].size(); i++) {
            Edge& e = g[u][i];
            if (e.cap > 0 && level[u] < level[e.to]) {
                long long d = dfs(e.to, t, min(f, e.cap));
                if (d > 0) { e.cap -= d; g[e.to][e.rev].cap += d; return d; }
            }
        }
        return 0;
    }
    long long maxflow(int s, int t) {
        long long flow = 0;
        while (bfs(s, t)) {
            fill(it.begin(), it.end(), 0);
            long long f;
            while ((f = dfs(s, t, LLONG_MAX)) > 0) flow += f;
        }
        return flow;
    }
};

int main() {
    int n, m;
    cin >> n >> m;
    Dinic dinic(n + 1);
    for (int i = 0; i < m; i++) {
        int a, b; long long c; cin >> a >> b >> c;
        dinic.addEdge(a, b, c);
    }
    cout << dinic.maxflow(1, n) << "\n";
    return 0;
}
\end{lstlisting}

\begin{complexitybox}
\textbf{Time Complexity.} $O(V^2 E)$ worst case (Dinic), far faster in
practice. \textbf{Space Complexity.} $O(V+E)$.
\end{complexitybox}

\begin{takeawaybox}
Maximum throughput is maximum flow; Dinic's level-graph blocking flow is the
standard efficient solver. This one algorithm is the engine behind the next
three problems -- min cut, matching, and edge-disjoint paths -- each just a
different network.
\end{takeawaybox}

\section{Police Chase}
\label{sec:cses-police-chase}

\problemstatement{Given an undirected graph of $n$ intersections and $m$
streets, find the minimum number of streets to remove so that node $1$ is
disconnected from node $n$, and output those streets.}

\begin{patternbox}
``Fewest edges to disconnect $s$ from $t$'' is the \textbf{minimum cut},
which by the max-flow--min-cut theorem equals the \textbf{maximum flow} with
every edge given capacity $1$. The cut edges are those crossing from the
source-reachable side to the rest after the flow saturates.
\end{patternbox}

\subsection*{Max flow, then read the cut}

Model each street as an undirected unit-capacity edge (a pair of directed
edges). The max flow from $1$ to $n$ equals the minimum number of streets
whose removal disconnects them. After computing it, BFS in the residual
graph from node $1$: the reachable set $S$ is one side of the cut. Every
original street with exactly one endpoint in $S$ is a cut edge -- print
those. Their count equals the max-flow value.

\begin{ideabox}[The Reachable Side Defines the Cut]
When no augmenting path remains, the nodes still reachable from the source
in the residual graph are separated from the rest by saturated edges. The
original edges straddling that boundary are precisely a minimum edge cut.
\end{ideabox}

\subsection*{Algorithm}

\begin{enumerate}
  \item Add each street as a unit-capacity edge in both directions; run max
        flow $1\to n$.
  \item BFS the residual graph from $1$ to get the reachable set $S$.
  \item Output every original street with one endpoint in $S$ and one
        outside.
\end{enumerate}

\subsection*{Dry run}

\dryrun{Path $1\!-\!2\!-\!3$ with $n=3$; a single street between $2$ and
$3$ (and $1,2$).}

\begin{itemize}
  \item Max flow $=1$; residual-reachable set from $1$ is $\{1,2\}$.
  \item The street $2\!-\!3$ crosses the boundary -- remove it (one
        street).
\end{itemize}

\subsection*{C++ implementation}

\begin{lstlisting}
#include <bits/stdc++.h>
using namespace std;

struct Dinic {
    struct Edge { int to; long long cap; int rev; };
    vector<vector<Edge>> g;
    vector<int> level, it;
    Dinic(int n) : g(n), level(n), it(n) {}
    void addEdge(int u, int v, long long c) {
        g[u].push_back({v, c, (int)g[v].size()});
        g[v].push_back({u, c, (int)g[u].size() - 1});   // undirected: both cap c
    }
    bool bfs(int s, int t) {
        fill(level.begin(), level.end(), -1);
        queue<int> q; level[s] = 0; q.push(s);
        while (!q.empty()) {
            int u = q.front(); q.pop();
            for (auto& e : g[u])
                if (e.cap > 0 && level[e.to] < 0) { level[e.to] = level[u] + 1; q.push(e.to); }
        }
        return level[t] >= 0;
    }
    long long dfs(int u, int t, long long f) {
        if (u == t) return f;
        for (int& i = it[u]; i < (int)g[u].size(); i++) {
            Edge& e = g[u][i];
            if (e.cap > 0 && level[u] < level[e.to]) {
                long long d = dfs(e.to, t, min(f, e.cap));
                if (d > 0) { e.cap -= d; g[e.to][e.rev].cap += d; return d; }
            }
        }
        return 0;
    }
    long long maxflow(int s, int t) {
        long long flow = 0;
        while (bfs(s, t)) {
            fill(it.begin(), it.end(), 0);
            long long f;
            while ((f = dfs(s, t, LLONG_MAX)) > 0) flow += f;
        }
        return flow;
    }
};

int main() {
    int n, m;
    cin >> n >> m;
    Dinic dinic(n + 1);
    vector<pair<int,int>> streets(m);
    for (auto& [a, b] : streets) { cin >> a >> b; dinic.addEdge(a, b, 1); }

    long long cut = dinic.maxflow(1, n);
    // residual reachability from source
    vector<bool> reach(n + 1, false);
    queue<int> q; q.push(1); reach[1] = true;
    while (!q.empty()) {
        int u = q.front(); q.pop();
        for (auto& e : dinic.g[u])
            if (e.cap > 0 && !reach[e.to]) { reach[e.to] = true; q.push(e.to); }
    }

    cout << cut << "\n";
    for (auto& [a, b] : streets)
        if (reach[a] != reach[b]) cout << a << " " << b << "\n";
    return 0;
}
\end{lstlisting}

\begin{complexitybox}
\textbf{Time Complexity.} $O(V^2E)$ for the flow; the cut extraction is
$O(V+E)$. \textbf{Space Complexity.} $O(V+E)$.
\end{complexitybox}

\begin{takeawaybox}
Minimum edge cut equals maximum flow with unit capacities, and the cut edges
straddle the residual-reachable frontier. Max-flow--min-cut is one of the
most powerful reductions in graph theory -- ``disconnect cheaply'' becomes
``push maximally.''
\end{takeawaybox}

\section{School Dance}
\label{sec:cses-school-dance}

\problemstatement{There are $n$ boys and $m$ girls, and $k$ acceptable
boy--girl pairs. Form the maximum number of dancing pairs so that each
person is in at most one pair, and output the pairs.}

\begin{patternbox}
Maximum \textbf{bipartite matching} is a \textbf{maximum flow}: add a super-
source into every boy, edges for each acceptable pair, and every girl into a
super-sink, all with capacity $1$. The max flow equals the maximum matching,
and the used pair-edges are the matching.
\end{patternbox}

\subsection*{Matching as unit-capacity flow}

Build nodes: source $S$, boys $1\ldots n$, girls $1\ldots m$, sink $T$. Add
$S\to\text{boy}$ (cap $1$, so each boy is used at most once),
$\text{boy}\to\text{girl}$ for each acceptable pair (cap $1$), and
$\text{girl}\to T$ (cap $1$). A unit of flow from $S$ to $T$ traces a
boy--girl pairing, and capacity $1$ everywhere enforces ``at most one
partner.'' The max flow is the largest matching; a boy--girl edge that ended
up saturated is a matched pair.

\begin{ideabox}[Unit Capacities Enforce ``At Most One'']
Capacity-$1$ edges out of the source (per boy) and into the sink (per girl)
guarantee each person is matched once. Integral max flow then corresponds
exactly to a maximum matching -- the classic flow reduction.
\end{ideabox}

\subsection*{Algorithm}

\begin{enumerate}
  \item Build the source/boys/girls/sink network with unit capacities.
  \item Run max flow; its value is the matching size.
  \item For each pair-edge boy$\to$girl whose forward capacity is now $0$,
        output the pair.
\end{enumerate}

\subsection*{Dry run}

\dryrun{Boys $\{1,2\}$, girls $\{1,2\}$; pairs $(1,1),(1,2),(2,1)$.}

\begin{itemize}
  \item Match boy $1$--girl $2$ and boy $2$--girl $1$: two pairs.
  \item Max flow $=2$.
\end{itemize}

\subsection*{C++ implementation}

\begin{lstlisting}
#include <bits/stdc++.h>
using namespace std;

struct Dinic {
    struct Edge { int to; long long cap; int rev; };
    vector<vector<Edge>> g;
    vector<int> level, it;
    Dinic(int n) : g(n), level(n), it(n) {}
    int addEdge(int u, int v, long long c) {
        g[u].push_back({v, c, (int)g[v].size()});
        g[v].push_back({u, 0, (int)g[u].size() - 1});
        return (int)g[u].size() - 1;                 // index of forward edge
    }
    bool bfs(int s, int t) {
        fill(level.begin(), level.end(), -1);
        queue<int> q; level[s] = 0; q.push(s);
        while (!q.empty()) {
            int u = q.front(); q.pop();
            for (auto& e : g[u])
                if (e.cap > 0 && level[e.to] < 0) { level[e.to] = level[u] + 1; q.push(e.to); }
        }
        return level[t] >= 0;
    }
    long long dfs(int u, int t, long long f) {
        if (u == t) return f;
        for (int& i = it[u]; i < (int)g[u].size(); i++) {
            Edge& e = g[u][i];
            if (e.cap > 0 && level[u] < level[e.to]) {
                long long d = dfs(e.to, t, min(f, e.cap));
                if (d > 0) { e.cap -= d; g[e.to][e.rev].cap += d; return d; }
            }
        }
        return 0;
    }
    long long maxflow(int s, int t) {
        long long flow = 0;
        while (bfs(s, t)) {
            fill(it.begin(), it.end(), 0);
            long long f;
            while ((f = dfs(s, t, LLONG_MAX)) > 0) flow += f;
        }
        return flow;
    }
};

int main() {
    int n, m, k;
    cin >> n >> m >> k;
    int S = 0, T = n + m + 1;                        // boys 1..n, girls n+1..n+m
    Dinic dinic(T + 1);
    for (int i = 1; i <= n; i++) dinic.addEdge(S, i, 1);
    for (int j = 1; j <= m; j++) dinic.addEdge(n + j, T, 1);

    vector<array<int,3>> pairEdge;                   // {boy, girl, forward edge index}
    for (int i = 0; i < k; i++) {
        int a, b; cin >> a >> b;                     // boy a, girl b
        int idx = dinic.addEdge(a, n + b, 1);
        pairEdge.push_back({a, n + b, idx});
    }

    cout << dinic.maxflow(S, T) << "\n";
    for (auto& [boy, girlNode, idx] : pairEdge)
        if (dinic.g[boy][idx].cap == 0)              // saturated => matched
            cout << boy << " " << girlNode - n << "\n";
    return 0;
}
\end{lstlisting}

\begin{complexitybox}
\textbf{Time Complexity.} $O(E\sqrt{V})$ for bipartite matching via Dinic.
\textbf{Space Complexity.} $O(V+E)$.
\end{complexitybox}

\begin{takeawaybox}
Bipartite matching is unit-capacity max flow with a super-source and super-
sink; saturated pair-edges are the matches. This reduction makes matching,
assignment, and covering problems all instances of flow.
\end{takeawaybox}

\section{Distinct Routes}
\label{sec:cses-distinct-routes}

\problemstatement{Given a directed graph of $n$ nodes and $m$ edges, find
the maximum number of \emph{edge-disjoint} routes from node $1$ to node
$n$, and output each route. Every edge may be used by at most one route.}

\begin{patternbox}
The maximum number of edge-disjoint paths from $s$ to $t$ equals the
\textbf{maximum flow} when every edge has capacity $1$ (Menger's theorem).
After computing the flow, \textbf{decompose} it into that many paths by
following edges that carry flow.
\end{patternbox}

\subsection*{Max flow, then decompose into paths}

Give every edge capacity $1$ and compute the max flow from $1$ to $n$; its
value $f$ is the number of edge-disjoint routes. An original edge $u\to v$
carries flow iff its forward capacity dropped to $0$ (equivalently its
reverse now has capacity $1$). To extract the routes, walk from the source:
at each node follow one unused flow-carrying out-edge, mark it consumed, and
continue to the sink; repeat $f$ times. Because the flow is edge-disjoint,
these walks partition the flow edges into $f$ paths.

\begin{ideabox}[Flow Decomposition Yields the Paths]
A unit $s$--$t$ flow of value $f$ always splits into $f$ edge-disjoint
$s\to t$ paths (plus possibly some cycles, which unit capacities here avoid).
Greedily tracing flow edges from the source recovers each path exactly once.
\end{ideabox}

\subsection*{Algorithm}

\begin{enumerate}
  \item Add each edge with capacity $1$; run max flow $1\to n$; let $f$ be
        its value.
  \item Collect, per node, the out-edges that carry flow.
  \item Trace $f$ paths from $1$ to $n$, consuming each flow edge once; print
        the count and paths.
\end{enumerate}

\subsection*{Dry run}

\dryrun{$1\!\to\!2$, $2\!\to\!4$, $1\!\to\!3$, $3\!\to\!4$; $n=4$.}

\begin{itemize}
  \item Max flow $=2$: routes $1\to2\to4$ and $1\to3\to4$.
  \item Both routes are edge-disjoint.
\end{itemize}

\subsection*{C++ implementation}

\begin{lstlisting}
#include <bits/stdc++.h>
using namespace std;

struct Dinic {
    struct Edge { int to; long long cap; int rev; };
    vector<vector<Edge>> g;
    vector<int> level, it;
    Dinic(int n) : g(n), level(n), it(n) {}
    void addEdge(int u, int v, long long c) {
        g[u].push_back({v, c, (int)g[v].size()});
        g[v].push_back({u, 0, (int)g[u].size() - 1});
    }
    bool bfs(int s, int t) {
        fill(level.begin(), level.end(), -1);
        queue<int> q; level[s] = 0; q.push(s);
        while (!q.empty()) {
            int u = q.front(); q.pop();
            for (auto& e : g[u])
                if (e.cap > 0 && level[e.to] < 0) { level[e.to] = level[u] + 1; q.push(e.to); }
        }
        return level[t] >= 0;
    }
    long long dfs(int u, int t, long long f) {
        if (u == t) return f;
        for (int& i = it[u]; i < (int)g[u].size(); i++) {
            Edge& e = g[u][i];
            if (e.cap > 0 && level[u] < level[e.to]) {
                long long d = dfs(e.to, t, min(f, e.cap));
                if (d > 0) { e.cap -= d; g[e.to][e.rev].cap += d; return d; }
            }
        }
        return 0;
    }
    long long maxflow(int s, int t) {
        long long flow = 0;
        while (bfs(s, t)) {
            fill(it.begin(), it.end(), 0);
            long long f;
            while ((f = dfs(s, t, LLONG_MAX)) > 0) flow += f;
        }
        return flow;
    }
};

int main() {
    int n, m;
    cin >> n >> m;
    Dinic dinic(n + 1);
    for (int i = 0; i < m; i++) {
        int a, b; cin >> a >> b;
        dinic.addEdge(a, b, 1);
    }

    long long f = dinic.maxflow(1, n);
    cout << f << "\n";

    // flow-carrying out-edges: forward edge with cap now 0 (capacity was 1)
    vector<int> ptr(n + 1, 0);
    for (long long r = 0; r < f; r++) {
        vector<int> path;
        int u = 1;
        path.push_back(1);
        while (u != n) {
            auto& es = dinic.g[u];
            while (ptr[u] < (int)es.size() &&
                   !(es[ptr[u]].to != 0 && es[ptr[u]].cap == 0 &&
                     dinic.g[es[ptr[u]].to][es[ptr[u]].rev].cap > 0))
                ptr[u]++;
            auto& e = es[ptr[u]];
            dinic.g[e.to][e.rev].cap = 0;             // consume this flow edge
            u = e.to;
            path.push_back(u);
        }
        cout << path.size() << "\n";
        for (int v : path) cout << v << " ";
        cout << "\n";
    }
    return 0;
}
\end{lstlisting}

\begin{complexitybox}
\textbf{Time Complexity.} $O(V^2E)$ for the flow; decomposition is $O(V+E)$
across all paths. \textbf{Space Complexity.} $O(V+E)$.
\end{complexitybox}

\begin{takeawaybox}
The most edge-disjoint $s$--$t$ paths is unit-capacity max flow (Menger's
theorem), and flow decomposition recovers the actual routes. This closes the
flow chapter: throughput, cut, matching, and disjoint paths are one
algorithm applied to four networks.
\end{takeawaybox}

